\chapter{Complete Python signatures and native CLI options}\label{ch:api-signatures}

\section{Using this reference}

The signatures below are generated from the public Python implementation. Property entries have no call parentheses. Private helpers and native-handle constructors are implementation details and are excluded. Atom's dataclass constructor is described explicitly. The workflow chapters explain physical interpretation, examples, and validation. Python indices are zero-based unless a function explicitly states otherwise.

\begin{lstlisting}[language=Python]

Atom(symbol, x, y, z, r=1.0, g=1.0, b=1.0)

\end{lstlisting}

Coordinates are Cartesian Angstrom. r/\allowbreak{}g/\allowbreak{}b here are colour components, not atomic radii. Structure.atoms is an editable list; Structure.cell is three lattice-vector rows or None. Prefer add\_\allowbreak{}atom for automatic element colours. Structure translation/\allowbreak{}scaling edits in place; repeat and filter\_\allowbreak{}species return new structures. Grid operations return new data.

\section{Atom and Structure}

\subsection{Atom}

A single atom with Cartesian coordinates (Angstrom) and display colour.

\begin{samepage}

\begin{lstlisting}[language=Python]
Atom.copy()
\end{lstlisting}

Return a separate atom with the same coordinates and colour.

\end{samepage}



\subsection{Structure}

A collection of atoms with an optional periodic unit cell. atoms is a list of Atom objects with Cartesian positions in Angstrom. cell is three lattice-vector rows [[a1,a2,a3],[b1,b2,b3],[c1,c2,c3]], in Angstrom, or None for a nonperiodic cluster.

\begin{samepage}

\begin{lstlisting}[language=Python]
Structure.__init__()
\end{lstlisting}

Create an empty structure with atoms=[] and cell=None.

\end{samepage}



\begin{samepage}

\begin{lstlisting}[language=Python]
Structure.copy()
\end{lstlisting}

Return a separate structure with copied atoms and cell vectors.

\end{samepage}



\begin{samepage}

\begin{lstlisting}[language=Python]
Structure.add_atom(symbol: str, x: float, y: float, z: float, r: float=-1.0, g: float=-1.0, b: float=-1.0)
\end{lstlisting}

Append an atom and return it.  Default colour is the element's CPK colour.

\end{samepage}



\begin{samepage}

\begin{lstlisting}[language=Python]
Structure.remove_atom(index: int)
\end{lstlisting}

Remove and return the atom at index.

\end{samepage}



\begin{samepage}

\begin{lstlisting}[language=Python]
Structure.translate(dx: float, dy: float, dz: float)
\end{lstlisting}

Shift all atoms in-place by (dx, dy, dz) Angstrom.  Returns self.

\end{samepage}



\begin{samepage}

\begin{lstlisting}[language=Python]
Structure.scale(factor: float)
\end{lstlisting}

Scale all Cartesian coordinates and cell vectors by factor.  Returns self.

\end{samepage}



\begin{samepage}

\begin{lstlisting}[language=Python]
Structure.repeat(nx: int, ny: int=1, nz: int=1)
\end{lstlisting}

Return a new supercell replicated nx x ny x nz times.  Requires a unit cell.

\end{samepage}



\begin{samepage}

\begin{lstlisting}[language=Python]
Structure.filter_species(*symbols: str)
\end{lstlisting}

Return a new Structure containing only atoms of the given element symbols.

\end{samepage}



\begin{samepage}

\begin{lstlisting}[language=Python]
Structure.set_cell(a: float, b: float, c: float, alpha: float=90.0, beta: float=90.0, gamma: float=90.0)
\end{lstlisting}

Set the unit cell from lattice parameters (lengths in Angstrom, angles in degrees).
Uses the standard convention: a along x, b in the xy-plane.
Returns self.

\end{samepage}



\begin{samepage}

\begin{lstlisting}[language=Python]
Structure.view()
\end{lstlisting}

Open this structure in the AtomForge GUI (non-blocking).

\end{samepage}



\begin{samepage}

\begin{lstlisting}[language=Python]
Structure.view_notebook(width: int=620, height: int=460)
\end{lstlisting}

Explicitly display the inline 3D viewer in a Jupyter notebook cell.

\end{samepage}



\begin{samepage}

\begin{lstlisting}[language=Python]
Structure.save(path: str)
\end{lstlisting}

Save to file; format is inferred from the file extension.

\end{samepage}



\section{Structure file functions}

\begin{samepage}

\begin{lstlisting}[language=Python]
load(path: str)
\end{lstlisting}

Load a structure; format is inferred from the filename.

\end{samepage}



\begin{samepage}

\begin{lstlisting}[language=Python]
save(s: 'Structure', path: str)
\end{lstlisting}

Save a structure; format is inferred from the filename.

\end{samepage}



\section{Desktop viewer function}

\begin{samepage}

\begin{lstlisting}[language=Python]
view(s: 'Structure')
\end{lstlisting}

Open s in the AtomForge GUI (non-blocking).

The structure is serialised to a temporary extXYZ file and AtomForge is
launched as a detached subprocess.  The temp file is cleaned up after
AtomForge exits; a daemon cleanup thread allows this function to remain
non-blocking.

\end{samepage}



\section{Native Python builders}

\subsection{BuildResult}

Built structure and native diagnostics; output is a saved path or None.

Inspect stdout and stderr for algorithm warnings. Explicit output files and
native orientation sidecars are retained. Structure holds coordinates and
cell data; it does not import the sidecar's grain metadata.

\begin{samepage}

\begin{lstlisting}[language=Python]
builder_help(mode, *, executable=None)
\end{lstlisting}

Return installed native help for a registered builder mode.

\end{samepage}



\begin{samepage}

\begin{lstlisting}[language=Python]
build(mode, options=(), *, source=None, output=None, executable=None, timeout=300)
\end{lstlisting}

Run a native builder and return BuildResult without opening the GUI.

mode is bulk, gb, poly, nano, amorphous, sss, dislocation, or custom.
options is an argument sequence, e.g. ['--a', '3.61', '--atom', 'Cu 0 0 0'].
Repeated flags are retained; multi-component values occupy one argument.
source accepts a Structure or an input path. output optionally retains a
CIF, XYZ/\allowbreak{}extXYZ, VASP, PDB, or LAMMPS result and native sidecars; otherwise
temporary VASP files are cleaned up after loading. Coordinates are Angstrom.
Native failures raise subprocess.CalledProcessError with stdout/\allowbreak{}stderr;
timeout raises subprocess.TimeoutExpired. No command shell is used.

\end{samepage}



\begin{samepage}

\begin{lstlisting}[language=Python]
wulff(source, facets, *, radius=20, vacuum=5, output=None, executable=None, timeout=300)
\end{lstlisting}

Build a Wulff particle using (h, k, l, energy) facet families.

Positive relative surface energies set plane distances; radius is the
maximum plane distance in Angstrom, not the farthest vertex radius.
Uses symmetry expansion and the same native engine as the desktop.
Returns BuildResult. Requires an executable supporting --shape wulff.

\end{samepage}



\section{Electronic data types and operations}

\subsection{Grid}

Scalar grid in Angstrom; values use x-fastest order (z, y, x loops).

cell contains three lattice/\allowbreak{}span vectors, as in Structure.cell.
Periodic grids omit endpoints; finite grids include them. Operations return
new grids. unit='e/\allowbreak{}A\textasciicircum{}3' enables density-derived energy calculations.

\begin{samepage}

\begin{lstlisting}[language=Python]
Grid.__init__(shape, cell, values, origin=(0, 0, 0), periodic=False, unit='raw')
\end{lstlisting}

Create a scalar grid. shape contains three integers >=2; values has exactly their product in x-fastest order. cell has three vectors in Angstrom. origin is Cartesian. Periodic grids omit duplicate endpoints; finite grids include endpoints. The grid is limited to 100 million samples.

\end{samepage}



\begin{samepage}

\begin{lstlisting}[language=Python]
Grid.shape
\end{lstlisting}

Read-only (nx, ny, nz) tuple.

\end{samepage}



\begin{samepage}

\begin{lstlisting}[language=Python]
Grid.cell
\end{lstlisting}

Read-only tuple of the three span/\allowbreak{}lattice vectors in Angstrom.

\end{samepage}



\begin{samepage}

\begin{lstlisting}[language=Python]
Grid.origin
\end{lstlisting}

Read-only Cartesian grid origin in Angstrom.

\end{samepage}



\begin{samepage}

\begin{lstlisting}[language=Python]
Grid.periodic
\end{lstlisting}

Read-only Boolean periodic-grid convention.

\end{samepage}



\begin{samepage}

\begin{lstlisting}[language=Python]
Grid.values
\end{lstlisting}

Return an independent list of scalar samples.

\end{samepage}



\begin{samepage}

\begin{lstlisting}[language=Python]
Grid.unit
\end{lstlisting}

Read-only scalar unit label; do not infer it from a plotted colour.

\end{samepage}



\begin{samepage}

\begin{lstlisting}[language=Python]
Grid.name
\end{lstlisting}

Read-only field name, including imported/\allowbreak{}derived channel identification.

\end{samepage}



\begin{samepage}

\begin{lstlisting}[language=Python]
Grid.sites
\end{lstlisting}

Return (atomic\_\allowbreak{}number, x, y, z) rows; unknown species have number 0.

\end{samepage}



\begin{samepage}

\begin{lstlisting}[language=Python]
Grid.with_sites(sites)
\end{lstlisting}

Return a grid with supplied (atomic\_\allowbreak{}number, x, y, z) site rows; positions are Cartesian Angstrom.

\end{samepage}



\begin{samepage}

\begin{lstlisting}[language=Python]
Grid.__add__(other)
\end{lstlisting}

Return pointwise addition of aligned grids.

\end{samepage}



\begin{samepage}

\begin{lstlisting}[language=Python]
Grid.__sub__(other)
\end{lstlisting}

Return pointwise subtraction of aligned grids, self minus other.

\end{samepage}



\begin{samepage}

\begin{lstlisting}[language=Python]
Grid.__mul__(other)
\end{lstlisting}

Return pointwise grid multiplication or scalar scaling.

\end{samepage}



\begin{samepage}

\begin{lstlisting}[language=Python]
Grid.__truediv__(other)
\end{lstlisting}

Return pointwise grid division or scalar division. Scalar zero is rejected.

\end{samepage}



\begin{samepage}

\begin{lstlisting}[language=Python]
Grid.resample(target)
\end{lstlisting}

Interpolate this field on target grid geometry. This does not register displaced atoms or correct incompatible units.

\end{samepage}



\begin{samepage}

\begin{lstlisting}[language=Python]
Grid.density_difference(*references, weights=None)
\end{lstlisting}

Return self - sum(weight * reference), in e/\allowbreak{}A\textasciicircum{}3.

Use fragment densities calculated in the same cell and geometry.
Alignment is required; resample explicitly when appropriate.

\end{samepage}



\begin{samepage}

\begin{lstlisting}[language=Python]
Grid.threshold_mask(low, high)
\end{lstlisting}

Binary mask for the inclusive interval [low, high].

\end{samepage}



\begin{samepage}

\begin{lstlisting}[language=Python]
Grid.boolean(other, operation='intersection')
\end{lstlisting}

Combine binary masks: union, intersection, difference (self-other), xor.

\end{samepage}



\begin{samepage}

\begin{lstlisting}[language=Python]
Grid.invert_mask()
\end{lstlisting}

Complement a binary mask within its cell.

\end{samepage}



\begin{samepage}

\begin{lstlisting}[language=Python]
Grid.apply_mask(mask)
\end{lstlisting}

Keep density inside a binary mask, preserving the density unit.

\end{samepage}



\begin{samepage}

\begin{lstlisting}[language=Python]
Grid.split_density()
\end{lstlisting}

Return accumulation and positive depletion-magnitude fields.

\end{samepage}



\begin{samepage}

\begin{lstlisting}[language=Python]
Grid.charge_summary()
\end{lstlisting}

Electron accumulation, depletion magnitude and net change.

Apply to a difference density. These integrals describe redistribution;
they do not assign charge transfer to individual atoms or fragments.

\end{samepage}



\begin{samepage}

\begin{lstlisting}[language=Python]
Grid.cumulative_charge(axis=2)
\end{lstlisting}

(Normal distance in A, cumulative electrons) from the lower cell face.

The final row includes the full cell. Periodic profiles depend on the
chosen cell origin; finite grids use trapezoidal quadrature.

\end{samepage}



\begin{samepage}

\begin{lstlisting}[language=Python]
Grid.as_periodic(tolerance=1e-08)
\end{lstlisting}

Drop matching endpoint planes after checking all three boundaries.

Only use when the finite grid spans a complete periodic cell. This
conversion checks values, but cannot establish physical periodicity.

\end{samepage}



\begin{samepage}

\begin{lstlisting}[language=Python]
Grid.gradient()
\end{lstlisting}

Return Cartesian x/\allowbreak{}y/\allowbreak{}z derivatives, including triclinic metric.

\end{samepage}



\begin{samepage}

\begin{lstlisting}[language=Python]
Grid.laplacian()
\end{lstlisting}

Return the Cartesian Laplacian grid, including cell geometry; units are field units per Angstrom squared.

\end{samepage}



\begin{samepage}

\begin{lstlisting}[language=Python]
Grid.energy_density(floor=1e-12)
\end{lstlisting}

Return approximate kinetic, potential and total density in eV/\allowbreak{}A\textasciicircum{}3.

Nonnegative electron density is required; samples <= floor (e/\allowbreak{}A\textasciicircum{}3)
are masked to zero. These are gradient-expansion/\allowbreak{}local-virial fields,
not the kinetic energy evaluated from Kohn-Sham orbitals.

\end{samepage}



\begin{samepage}

\begin{lstlisting}[language=Python]
Grid.smooth(sigma, radius=2)
\end{lstlisting}

Return Gaussian-smoothed grid. sigma is in Angstrom; radius bounds the kernel in grid steps.

\end{samepage}



\begin{samepage}

\begin{lstlisting}[language=Python]
Grid.line_profile(start, end, count=200)
\end{lstlisting}

Return distance/\allowbreak{}value rows on the Cartesian segment start to end (Angstrom), using count samples.

\end{samepage}



\begin{samepage}

\begin{lstlisting}[language=Python]
Grid.planar_average(axis=2)
\end{lstlisting}

Return normal distance in Angstrom and area-averaged field.

\end{samepage}



\begin{samepage}

\begin{lstlisting}[language=Python]
Grid.macroscopic_average(axis=2, window=3)
\end{lstlisting}

Centered box average over an odd number of planar samples.

\end{samepage}



\begin{samepage}

\begin{lstlisting}[language=Python]
Grid.section(origin, u, v, shape=(100, 100))
\end{lstlisting}

Sample origin + su + tv for 0<=s,t<=1 (Cartesian Angstrom).

Returns a two-layer grid with identical layers; the first layer is
the section. Use section\_\allowbreak{}values/\allowbreak{}contours for 2D data, not its volume.

\end{samepage}



\begin{samepage}

\begin{lstlisting}[language=Python]
Grid.section_values
\end{lstlisting}

Return the first z layer as rows indexed by y, with x-fastest entries. Use on a section result.

\end{samepage}



\begin{samepage}

\begin{lstlisting}[language=Python]
Grid.contours(level)
\end{lstlisting}

Return endpoint pairs for contour segments on grid z=0.

\end{samepage}



\begin{samepage}

\begin{lstlisting}[language=Python]
Grid.integrate()
\end{lstlisting}

Return the whole-grid scalar integral with volume quadrature; density gives electrons.

\end{samepage}



\begin{samepage}

\begin{lstlisting}[language=Python]
Grid.integrate_sphere(center, radius)
\end{lstlisting}

Return the scalar integral within a sphere with Cartesian center and radius in Angstrom.

\end{samepage}



\begin{samepage}

\begin{lstlisting}[language=Python]
Grid.voronoi_integrate(sites=None)
\end{lstlisting}

Return (integral, volume) per site; equidistant samples share weight.

\end{samepage}



\begin{samepage}

\begin{lstlisting}[language=Python]
Grid.peaks(limit=0)
\end{lstlisting}

Grid-local maxima as (x,y,z,value); deterministic adjacent tie break.

\end{samepage}



\begin{samepage}

\begin{lstlisting}[language=Python]
Grid.structure_factors(indices)
\end{lstlisting}

Return (h,k,l,real,imag); phase convention is exp(+2pii*h.r).

\end{samepage}



\begin{samepage}

\begin{lstlisting}[language=Python]
Grid.fourier_synthesis(reflections)
\end{lstlisting}

Supply complete (h,k,l,real,imag) rows, including conjugate pairs.

\end{samepage}



\begin{samepage}

\begin{lstlisting}[language=Python]
Grid.patterson()
\end{lstlisting}

Return the periodic autocorrelation/\allowbreak{}Patterson field. See the operation reference for normalisation and interpretation.

\end{samepage}



\begin{samepage}

\begin{lstlisting}[language=Python]
Grid.atomic_structure_factors(indices, scattering, occupancies=None, b_factors=None)
\end{lstlisting}

Explicit complex scattering[reflection][site], occupancy and B(A\textasciicircum{}2).

\end{samepage}



\begin{samepage}

\begin{lstlisting}[language=Python]
Grid.ewald(charges, alpha, real_cutoff, reciprocal_cutoff)
\end{lstlisting}

Neutral periodic point charges; return energy(eV), potentials(V).

alpha is in 1/\allowbreak{}A, real\_\allowbreak{}cutoff in A and reciprocal\_\allowbreak{}cutoff in 1/\allowbreak{}A
(wavevector including 2*pi). Check convergence by increasing cutoffs.

\end{samepage}



\begin{samepage}

\begin{lstlisting}[language=Python]
Grid.isosurface(level, color=None)
\end{lstlisting}

Return a Surface at level; optional aligned color grid supplies per-vertex scalar colour values.

\end{samepage}



\begin{samepage}

\begin{lstlisting}[language=Python]
Grid.save(path, format=None)
\end{lstlisting}

Write this field. Explicit format is xsf, cube, or vasp. Without format, .cube/\allowbreak{}.cub and .xsf are recognised; other suffixes default to VASP.

\end{samepage}



\subsection{Volume}

An imported collection of scalar channels and their atomic sites.

\begin{samepage}

\begin{lstlisting}[language=Python]
Volume.fields
\end{lstlisting}

Read-only tuple of Grid channels. Select with a zero-based index after checking names and units.

\end{samepage}



\begin{samepage}

\begin{lstlisting}[language=Python]
Volume.save(path, format)
\end{lstlisting}

Write the imported volume using an explicit supported format. Check the selected format's channel capabilities.

\end{samepage}



\subsection{Surface}

Triangles with per-vertex scalar coloring; vertices are not welded.

\begin{samepage}

\begin{lstlisting}[language=Python]
Surface.vertices
\end{lstlisting}

Return the native triangle vertex rows with scalar data; triangles are not welded into a shared-vertex topology.

\end{samepage}



\begin{samepage}

\begin{lstlisting}[language=Python]
Surface.save(path)
\end{lstlisting}

Write OBJ (scalar comments) or .ply (a scalar vertex property).

\end{samepage}



\begin{samepage}

\begin{lstlisting}[language=Python]
load_volume(path, quantity='auto', cube_coordinates='bohr')
\end{lstlisting}

Read VASP, Cube or XSF. Cube/\allowbreak{}XSF 'auto' keeps scalar units raw.

quantity: density, potential, elf, raw, auto. Cube density/\allowbreak{}potential are
converted from atomic units. XSF scalars are assumed e/\allowbreak{}A\textasciicircum{}3 or eV only
when explicitly requested; producer-specific units must be converted.

\end{samepage}



\section{Electronic model helpers}

\begin{samepage}

\begin{lstlisting}[language=Python]
reciprocal_vectors(cell)
\end{lstlisting}

Reciprocal vectors without 2*pi, in inverse Angstrom.

\end{samepage}



\begin{samepage}

\begin{lstlisting}[language=Python]
scattering_factor(s, a, b, c=0.0)
\end{lstlisting}

Gaussian X-ray form factor sum(a\_\allowbreak{}iexp(-b\_\allowbreak{}is\textasciicircum{}2))+c.

Supply published coefficients and respect their stated valid s=sin(theta)/\allowbreak{}
wavelength range. No anomalous-dispersion correction is added implicitly.

\end{samepage}



\begin{samepage}

\begin{lstlisting}[language=Python]
model_density(geometry, coefficients=None, scattering_lengths=None, occupancies=None, b_factors=None)
\end{lstlisting}

Fourier model electron or nuclear scattering-length density.

Exactly one mapping keyed by atomic number is required: coefficients
\{Z: (a, b, c)\} for X-rays or scattering\_\allowbreak{}lengths \{Z: length\} for
neutrons. Length units are caller-defined. Even-grid Nyquist planes are
omitted to retain a real, conjugate-symmetric finite Fourier expansion.
These are free-atom scattering models, not self-consistent densities.

\end{samepage}



\begin{samepage}

\begin{lstlisting}[language=Python]
miller_section(grid, hkl, offset=0.0, width=5.0, height=5.0, shape=(100, 100))
\end{lstlisting}

Section centered on hx+ky+l*z=offset (fractional cell coordinates).

Width/\allowbreak{}height are Angstrom. Finite inputs must contain the whole section.
The orthonormal in-plane orientation is chosen deterministically.

\end{samepage}



\section{Structural analyses and sculpting}

\begin{samepage}

\begin{lstlisting}[language=Python]
analyze(source, method, *, executable=None, timeout=300, **parameters)
\end{lstlisting}

Run cna, rdf, adf, sro or interstitial and return named CSV rows.

Parameter names match CLI flags with underscores replacing hyphens.
no\_\allowbreak{}pbc=True and raw=True are switches. SRO currently uses the
native shell convention and does not accept a no\_\allowbreak{}pbc override.

\end{samepage}



\begin{samepage}

\begin{lstlisting}[language=Python]
cna(source, **parameters)
\end{lstlisting}

Common-neighbour indices and full-environment classification per atom.

\end{samepage}



\begin{samepage}

\begin{lstlisting}[language=Python]
rdf(source, **parameters)
\end{lstlisting}

Radial distribution, directed raw pair counts and coordination.

\end{samepage}



\begin{samepage}

\begin{lstlisting}[language=Python]
adf(source, **parameters)
\end{lstlisting}

Angular distribution of unordered neighbour pairs, in degrees.

\end{samepage}



\begin{samepage}

\begin{lstlisting}[language=Python]
sro(source, **parameters)
\end{lstlisting}

Warren-Cowley and Rao-Curtin short-range order by neighbour shell.

\end{samepage}



\begin{samepage}

\begin{lstlisting}[language=Python]
interstitial_sites(source, **parameters)
\end{lstlisting}

Interstitial candidates, clearance, volume and coordination.

\end{samepage}



\begin{samepage}

\begin{lstlisting}[language=Python]
sculpt(source, slabs, *, repeat=(1, 1, 1), periodic=False, executable=None, timeout=300)
\end{lstlisting}

Apply intersecting Miller slabs with the desktop sculptor engine.

Each slab is (h,k,l,lower,upper). Manual bounds are A relative to the
supercell centroid; periodic bounds mean integer start plane and period
count, using the native structural-period convention.

\end{samepage}



\section{Electronic recipes and batch jobs}

\begin{samepage}

\begin{lstlisting}[language=Python]
apply_operation(grid, operation, parameters=None, reference=None)
\end{lstlisting}

Evaluate one allowlisted operation; return a grid, fields, surface or table.

\end{samepage}



\begin{samepage}

\begin{lstlisting}[language=Python]
run_pipeline(grid, steps, *, base_directory='.')
\end{lstlisting}

Run JSON-compatible steps and retain all named intermediate results.

Each step contains operation, optional parameters, name,
input (an earlier grid name), and reference (earlier name or file).
field selects a component when an operation returns multiple fields.
File references resolve relative to base\_\allowbreak{}directory. Recipes cannot call
arbitrary Python functions or execute shell commands.

\end{samepage}



\begin{samepage}

\begin{lstlisting}[language=Python]
save_result(result, path, format=None)
\end{lstlisting}

Write a grid/\allowbreak{}mesh, numeric CSV, JSON summary, or numbered grid components.

\end{samepage}



\begin{samepage}

\begin{lstlisting}[language=Python]
batch_process(paths, steps, output_directory, *, quantity='auto', channel=0, format='xsf', base_directory='.', overwrite=False)
\end{lstlisting}

Apply one recipe to every input and write the last result plus provenance.

Refuses name collisions and existing outputs by default. Failed inputs are
reported individually; successful results remain available.

\end{samepage}



\section{Relaxation, dynamics and phonons}

\begin{samepage}

\begin{lstlisting}[language=Python]
to_ase(structure)
\end{lstlisting}

Copy species, Cartesian positions and cell into an ASE Atoms object.

\end{samepage}



\begin{samepage}

\begin{lstlisting}[language=Python]
from_ase(atoms)
\end{lstlisting}

Copy species, positions and a full-rank cell into an AtomForge Structure.

\end{samepage}



\subsection{RelaxationResult}



\begin{samepage}

\begin{lstlisting}[language=Python]
relax(structure, calculator, *, fmax=0.02, steps=200, relax_cell=False, trajectory=None)
\end{lstlisting}

Run BFGS with an explicitly supplied ASE calculator (including DFT engines).

converged must be checked: reaching the step limit is not convergence.
Calculator setup, pseudopotentials and electronic convergence are controlled
by the caller. The source Structure is never modified.

\end{samepage}



\begin{samepage}

\begin{lstlisting}[language=Python]
molecular_dynamics(structure, calculator, *, steps=100, timestep_fs=1, temperature_K=300, seed=0, sample_interval=1, velocities=None)
\end{lstlisting}

Run NVE velocity-Verlet; return frames and total energies in eV.

With velocities omitted, initialize a Maxwell-Boltzmann distribution using
seed. Supplied velocities use ASE's internal velocity units. Temperature is
an initialization target, not a thermostat. Inspect energy drift to select dt.

\end{samepage}



\begin{samepage}

\begin{lstlisting}[language=Python]
phonons(structure, calculator, qpoints, directory, *, supercell=(2, 2, 2), displacement=0.01)
\end{lstlisting}

Finite-displacement phonons; energies in eV and complex eigenvectors.

qpoints are fractional reciprocal coordinates. Negative energies denote
imaginary frequencies. A fresh cache directory prevents stale force reuse.

\end{samepage}



\begin{samepage}

\begin{lstlisting}[language=Python]
mode_frames(structure, mode, *, amplitude=0.2, frames=40)
\end{lstlisting}

Animate one Gamma-point mode, with maximum displacement amplitude in A.

\end{samepage}



\section{Bands and density of states}

\begin{samepage}

\begin{lstlisting}[language=Python]
read_bands(path)
\end{lstlisting}

Read EIGENVAL with validated dimensions; return kpoints, weights, spin energies/\allowbreak{}occupancies.

Energies are absolute VASP eigenvalues in eV. Fermi energy is not present in
EIGENVAL and must be supplied separately when plotting.

\end{samepage}



\begin{samepage}

\begin{lstlisting}[language=Python]
read_projections(path)
\end{lstlisting}

Read PROCAR, including spin and SOC projections via pymatgen.

Returns the native Procar object: data[spin][kpoint,band,ion,orbital],
eigenvalues, occupancies, orbital names and SOC xyz\_\allowbreak{}data when present.

\end{samepage}



\begin{samepage}

\begin{lstlisting}[language=Python]
read_dos(path, *, projected_labels=None)
\end{lstlisting}

Read DOSCAR total DOS and all projected columns without guessing LORBIT.

Total DOS columns are labelled by spin; projected columns require caller
labels or remain column\_\allowbreak{}1, column\_\allowbreak{}2, ... because DOSCAR does not identify
the orbital/\allowbreak{}SOC convention. Integrated total DOS is retained separately.

\end{samepage}



\begin{samepage}

\begin{lstlisting}[language=Python]
plot_bands(data, output, *, fermi_eV=0, distances=None, projections=None)
\end{lstlisting}

Plot bands or fat bands; projections map spin to nonnegative (k,band) weights.

Without reciprocal-space distances the horizontal axis is k-point index,
avoiding an incorrect metric for non-orthogonal reciprocal cells.

\end{samepage}



\begin{samepage}

\begin{lstlisting}[language=Python]
plot_dos(data, output, *, sites=(), columns=None)
\end{lstlisting}

Save total and optionally site-projected DOS relative to the file's Fermi energy.

\end{samepage}



\section{Population analysis and orbital reconstruction}

\begin{samepage}

\begin{lstlisting}[language=Python]
hirshfeld(density, proatoms, *, reference_electrons=None)
\end{lstlisting}

Stockholder partition using supplied neutral proatom density grids.

All grids must share geometry, units and periodicity. The caller supplies
physically appropriate neutral-atom densities; this routine does not invent
Gaussian substitutes. Returns electron populations and optional net atomic
charges, reference\_\allowbreak{}electrons - population. A valence density requires valence
proatoms/\allowbreak{}reference electron counts; do not mix all-electron and valence data.

\end{samepage}



\begin{samepage}

\begin{lstlisting}[language=Python]
bader(path, *, reference=None, potcar=None, executable=None, format='vasp')
\end{lstlisting}

Run the Henkelman Bader executable through pymatgen.

reference can be AECCAR0+AECCAR2 for basin definition while integrating the
supplied valence density. POTCAR supplies reference valence counts for charge
transfer; raw electron populations are always returned. No Voronoi fallback.
Temporary solver inputs use six-decimal lattice/\allowbreak{}site coordinates for legacy
fixed-column readers; scalar samples retain 16 digits. Calls use pymatgen's
working-directory-based runner and should run in separate processes, not threads.

\end{samepage}



\begin{samepage}

\begin{lstlisting}[language=Python]
ddec(directory, *, atomic_densities=None, run=True)
\end{lstlisting}

Run/\allowbreak{}read Chargemol DDEC6 and CM5 results through pymatgen.

Requires the Chargemol executable and its atomic-density database. Input
directory contains the calculation's CHGCAR/\allowbreak{}POTCAR/\allowbreak{}AECCAR files. External
solver failures propagate; an approximate partition is never substituted.

\end{samepage}



\begin{samepage}

\begin{lstlisting}[language=Python]
orbital_density(wavecar, poscar, output, *, kpoint=0, band=0, spin=None, spinor=None, scale=2)
\end{lstlisting}

Export a selected WAVECAR orbital as PARCHG, loadable in the desktop.

Indices are zero-based. Uses pymatgen's plane-wave reconstruction; PAW
augmentation is not included, so this is not an all-electron density.
Occupation and k-point weights are not applied. A collinear spin selection
returns one component; omitted spin in ISPIN=2 returns total and difference.

\end{samepage}



\section{Trajectory conversion and animation}

\begin{samepage}

\begin{lstlisting}[language=Python]
read_trajectory(path, *, format=None, index=':')
\end{lstlisting}

Read XYZ/\allowbreak{}extXYZ, XDATCAR, LAMMPS dumps and other ASE-supported trajectories.

\end{samepage}



\begin{samepage}

\begin{lstlisting}[language=Python]
write_trajectory(frames, path, *, format=None)
\end{lstlisting}

Write a nonempty sequence of structures using an ASE trajectory format.

\end{samepage}



\begin{samepage}

\begin{lstlisting}[language=Python]
export_animation(frames, path, *, fps=12, rotation='20x,30y', radii=0.5)
\end{lstlisting}

Export a GIF with fixed spatial bounds and explicitly chosen atom radii.

\end{samepage}



\section{Diffraction}

\begin{samepage}

\begin{lstlisting}[language=Python]
powder_diffraction(structure, *, radiation='xray', wavelength=1.5406, two_theta=(0, 90))
\end{lstlisting}

Return pymatgen diffraction peaks; wavelength A, scattering angle degrees.

\end{samepage}



\begin{samepage}

\begin{lstlisting}[language=Python]
single_crystal(structure, indices, scattering)
\end{lstlisting}

Return hkl, reciprocal vectors (2*pi/\allowbreak{}A) and |F|\textasciicircum{}2 for explicit site factors.

scattering is [reflection][atom], permitting complex anomalous factors and
neutron isotope factors. Intensities omit instrument/\allowbreak{}polarization corrections.

\end{samepage}



\begin{samepage}

\begin{lstlisting}[language=Python]
plot_diffraction(pattern, output, *, fwhm=0.1, points=4000)
\end{lstlisting}

Plot Gaussian-broadened powder peaks; FWHM is in degrees of 2 theta.

\end{samepage}



\section{Reciprocal-space band grids}

\subsection{BandGrid}



\begin{samepage}

\begin{lstlisting}[language=Python]
BandGrid.surface(band, energy=None)
\end{lstlisting}

Triangulate E(k)=energy; the default is the file's Fermi energy.

\end{samepage}



\begin{samepage}

\begin{lstlisting}[language=Python]
read_bandgrid(path, *, energy_unit='eV', reciprocal_unit='1/A')
\end{lstlisting}

Read one three-dimensional BXSF block, preserving band labels.

BXSF does not declare universal physical units. Supply the producing
program's units explicitly when they differ from eV and inverse Angstrom.
Endpoint samples are retained, as required for general BXSF grids.

\end{samepage}



\begin{samepage}

\begin{lstlisting}[language=Python]
plot_surface(grid, band, output, *, energy=None)
\end{lstlisting}

Save a labeled reciprocal-space surface plot as PNG, SVG or PDF.

\end{samepage}



\section{Native builder options}

Use \code{AtomForge --help MODE} for detailed help. \code{AtomForge --build MODE --help} prints the general overview in this version. These option listings are extracted from the source help and checked against the normal executable. CLI defaults may differ from a dialog's remembered settings. The caret line continuations shown in help are for Windows cmd; use a single line or a subprocess argument list in PowerShell/Python.

\clearpage

\subsection{bulk}

\begin{lstlisting}
------------------------------------------------------------------
BULK CRYSTAL  (--build bulk)
------------------------------------------------------------------
  --system     <name>         Crystal system: triclinic | monoclinic |
                               orthorhombic | tetragonal | trigonal |
                               hexagonal | cubic   (default: cubic)
  --spacegroup <N>            Space-group number 1-230  (default: 225)
  --a   <Ang>                 Lattice parameter a       (default: 4.0)
  --b   <Ang>                 Lattice parameter b       (default: a)
  --c   <Ang>                 Lattice parameter c       (default: a)
  --alpha <deg>               Cell angle alpha          (default: 90)
  --beta  <deg>               Cell angle beta           (default: 90)
  --gamma <deg>               Cell angle gamma          (default: 90)
  --atom  "SYMBOL fx fy fz"   Asymmetric unit atom (fractional coords).
                               May be repeated for multiple atoms.
                               Example: --atom "Cu 0 0 0"
  --output <file>             Output file (format from extension)

Example:
  AtomForge --build bulk --system cubic --spacegroup 225 ^
            --a 3.61 --atom "Cu 0 0 0" --output cu_fcc.cif
\end{lstlisting}

\clearpage

\subsection{gb}

\begin{lstlisting}
------------------------------------------------------------------
CSL GRAIN BOUNDARY  (--build gb)
------------------------------------------------------------------
  --input  <file>             Reference structure (must have a unit cell)
  --axis   "u v w"            Rotation axis, e.g. "0 0 1"   (default: 0 0 1)
  --sigma  <N>                Exact sigma value to use
  --sigmamax <N>              Maximum sigma to search when --sigma is omitted
                               (picks the smallest available sigma, default: 100)
  --plane  <0|1|2>            GB-plane index within the sigma candidate (default: 0)
  --uca    <N>                Grain-A unit-cell repetitions along stack dir (default: 1)
  --ucb    <N>                Grain-B unit-cell repetitions along stack dir (default: 1)
  --vacuum <Ang>              Vacuum padding on each side  (default: 0)
  --gap    <Ang>              Interface gap between grains (default: 0)
  --overlap <Ang>             Overlap removal radius       (default: 0)
  --conventional              Keep conventional cell (skip primitive reduction)
  --output <file>             Output file (format from extension)

Example:
  AtomForge --build gb --input cu_fcc.cif --axis "0 0 1" ^
            --sigma 5 --plane 0 --uca 3 --ucb 3 ^
            --vacuum 5.0 --output cu_sigma5_gb.cif
\end{lstlisting}

\clearpage

\subsection{poly}

\begin{lstlisting}
------------------------------------------------------------------
POLYCRYSTAL  (--build poly)
------------------------------------------------------------------
  --input  <file>             Reference structure (must have a unit cell)
  --sizex  <Ang>              Box dimension X  (default: 50)
  --sizey  <Ang>              Box dimension Y  (default: 50)
  --sizez  <Ang>              Box dimension Z  (default: 50)
  --grains <N>                Number of Voronoi grains  (default: 8)
  --seed   <N>                Random seed for reproducibility  (default: 42)
  --euler  "phi1 Phi phi2"    Bunge Euler angles (deg) for grain 0, 1, ...
                               Specify once per grain in order. If fewer
                               than --grains values are given, remaining
                               grains get random orientations.
  --output <file>             Output file (format from extension)

Example:
  AtomForge --build poly --input cu_fcc.cif ^
            --sizex 100 --sizey 100 --sizez 100 ^
            --grains 12 --seed 7 --output cu_poly.cif
\end{lstlisting}

\clearpage

\subsection{nano}

\begin{lstlisting}
------------------------------------------------------------------
NANOCRYSTAL  (--build nano)
------------------------------------------------------------------
  --input  <file>             Reference crystal (must have a unit cell)
  --shape  <name>             sphere | ellipsoid | box | cylinder |
                               octahedron | truncated-octahedron |
                               cuboctahedron | wulff (default: sphere)
  --facet "h k l energy"      Wulff facet family; repeat for each family.
                               --radius sets maximum Wulff plane distance.
  --radius <Ang>              Sphere/octahedron/cuboctahedron radius (default: 15)
  --rx <Ang>                  Ellipsoid half-extent X  (default: 15)
  --ry <Ang>                  Ellipsoid half-extent Y  (default: 12)
  --rz <Ang>                  Ellipsoid half-extent Z  (default: 10)
  --hx <Ang>                  Box / truncated-oct half-extent X  (default: 15)
  --hy <Ang>                  Box / truncated-oct half-extent Y  (default: 15)
  --hz <Ang>                  Box / truncated-oct half-extent Z  (default: 15)
  --trunc <Ang>               Truncated-octahedron truncation radius  (default: 12)
  --cylradius <Ang>           Cylinder radius  (default: 12)
  --cylheight <Ang>           Cylinder height  (default: 30)
  --cylaxis   <0|1|2>         Cylinder axis: 0=X 1=Y 2=Z  (default: 2)
  --vacuum <Ang>              Vacuum padding around particle  (default: 5)
  --repa <N>                  Manual supercell replication A  (0 = auto)
  --repb <N>                  Manual supercell replication B  (0 = auto)
  --repc <N>                  Manual supercell replication C  (0 = auto)
  --output <file>             Output file (format from extension)

Example:
  AtomForge --build nano --input cu.cif --shape sphere --radius 20 ^
            --vacuum 5 --output cu_nano.xyz
\end{lstlisting}

\clearpage

\subsection{amorphous}

\begin{lstlisting}
------------------------------------------------------------------
AMORPHOUS STRUCTURE  (--build amorphous)
------------------------------------------------------------------
  --element "SYMBOL N"        Element species and atom count.
                               Repeat for each species in the mixture.
                               Example: --element "Si 80" --element "O 160"
  --density <g/cm3>           Target density for auto box size  (default: 2.0)
  --boxa <Ang>                Manual box length A (overrides --density)
  --boxb <Ang>                Manual box length B (overrides --density)
  --boxc <Ang>                Manual box length C (overrides --density)
  --scale <factor>            Cell scale factor before packing  (default: 1.0)
  --mindist "Z1 Z2 dist"      Minimum separation for element pair (Ang).
                               Repeat for multiple pairs.
  --covtol <frac>             Covalent-radii tolerance fraction  (default: 0.75)
  --seed   <N>                RNG seed (0 = time-based)  (default: 42)
  --attempts <N>              Max placement attempts per atom  (default: 1000)
  --output <file>             Output file (format from extension)

Example:
  AtomForge --build amorphous --element "Si 80" --element "O 160" ^
            --density 2.2 --seed 1 --output sio2.xyz
\end{lstlisting}

\clearpage

\subsection{sss}

\begin{lstlisting}
------------------------------------------------------------------
SUBSTITUTIONAL SOLID SOLUTION  (--build sss)
------------------------------------------------------------------
  --input  <file>             Host structure to substitute into
  --frac   SYM=frac,...       Comma-separated element=fraction pairs.
                               Fractions must sum to 1.
                               Example: --frac "Cu=0.7,Zn=0.3"
  --seed   <N>                RNG seed (0 = time-based)  (default: 12345)
  --output <file>             Output file (format from extension)

Example:
  AtomForge --build sss --input cu.cif ^'
            --frac "Cu=0.7,Zn=0.3" --seed 42 --output brass.cif
\end{lstlisting}

\clearpage

\subsection{dislocation}

\begin{lstlisting}
------------------------------------------------------------------
DISLOCATION  (--build dislocation)
------------------------------------------------------------------
  --input  <file>             Input structure (must have a unit cell)
  --output <file>             Output file (format from extension)
  --character <edge|screw|mixed>
                              Dislocation character (default: edge)
  --shape <halfplane|cylinder|sphere|ellipsoid|freeform>
                              Application region shape (default: halfplane)
  --manual-vectors            Disable automatic FCC/HCP/BCC vector presets
  --plane   "h k l"          Slip-plane Miller indices (default: 1 1 1)
  --burgers "u v w"          Burgers direction indices (default: 1 -1 0)
  --line    "u v w"          Line direction indices (default: 1 1 -2)
  --line-frac   "fx fy fz"   Dislocation line point in fractional coords
  --line-offset "x y z"      Cartesian offset added to line point (A)
  --bscale <f>                Burgers magnitude scale factor (default: 1.0)
  --bmag <A>                  Explicit Burgers magnitude in Angstroms (overrides auto-detect)
  --mixed-angle <deg>         Edge/screw mixing angle for mixed mode
  --nu <f>                    Poisson ratio (default: 0.33)
  --core <A>                  Core radius regularization (default: 1.2)
  --cutoff <A>                Radial attenuation distance, 0=off (default: 20)
  --line-half <A>             Half-length extent along line (default: 1e6)
  --cyl-radius <A>            Cylinder radius for shape=cylinder
  --sphere-radius <A>         Sphere radius for shape=sphere
  --ellipsoid "rx ry rz"     Ellipsoid radii for shape=ellipsoid
  --poly2d "x1 y1;x2 y2;..."  Freeform polygon for shape=freeform

Example:
  AtomForge --build dislocation --input base.cfg --character edge ^
            --manual-vectors --line "0 0 1" --burgers "1 0 0" ^
            --shape cylinder --cyl-radius 12 --output edge.cfg
\end{lstlisting}

\clearpage

\subsection{custom}

\begin{lstlisting}
------------------------------------------------------------------
CUSTOM MESH-FILL  (--build custom)
------------------------------------------------------------------
Fill a 3D mesh model with atoms from a reference crystal structure.

  --input   <file>            Reference crystal file (CIF, XYZ, VASP, PDB, ...)
  --mesh    <file>            3D model file (OBJ or STL)
  --scale   <factor>          Angstrom per model unit  (default: 1.0)
  --vacuum  <Ang>             Vacuum padding for output cell  (default: 5.0)
  --repa <N>                  Manual replication along a  (0 = auto)
  --repb <N>                  Manual replication along b  (0 = auto)
  --repc <N>                  Manual replication along c  (0 = auto)
  --rotx <deg>                Rotate crystal about X before fill  (default: 0)
  --roty <deg>                Rotate crystal about Y before fill  (default: 0)
  --rotz <deg>                Rotate crystal about Z before fill  (default: 0)
  --miller  "h k l"          Align crystal direction [hkl] with mesh +Z axis
                               (overrides --rotx/y/z when given)
  --output  <file>            Output file (format from extension)

Example:
  AtomForge --build custom --input cu.cif --mesh bunny.stl ^
            --scale 0.1 --vacuum 5 --output cu_bunny.xyz
\end{lstlisting}

\clearpage

\subsection{interface}

\begin{lstlisting}
------------------------------------------------------------------
HETEROGENEOUS INTERFACE  (--build interface)
------------------------------------------------------------------
  --layerA <file>             Structure for layer A (must have a unit cell)
  --layerB <file>             Structure for layer B (must have a unit cell)
  --nmax   <N>                Max supercell search index  (default: 4)
  --maxcells <N>              Max supercell area (unit cells)  (default: 16)
  --pick   <N>                Index of strain-matched supercell pair to use
                               (0 = best match, default: 0)
  --layersA <N>               Z repetitions of layer A  (default: 1)
  --layersB <N>               Z repetitions of layer B  (default: 1)
  --gap    <Ang>              Interface gap  (default: 2.0)
  --vacuum <Ang>              Vacuum above/below  (default: 10.0)
  --repx   <N>                XY repeat X  (default: 1)
  --repy   <N>                XY repeat Y  (default: 1)
  --output <file>             Output file (format from extension)

Example:
  AtomForge --build interface --layerA cu.cif --layerB ni.cif ^
            --nmax 4 --layersA 3 --layersB 3 ^
            --gap 2.0 --vacuum 10.0 --output cu_ni_interface.cif
\end{lstlisting}

\clearpage

\subsection{stacking-fault}

\begin{lstlisting}
STACKING FAULT (--build stacking-fault)
--input FILE --output FILE [--plane 0..6] [--layers 9] [--interval 0.1]
[--maximum 2] [--frame 0] [--orthogonal] [--sequence DIRECTORY]
Planes: 0 auto, 1 FCC111, 2 HCP basal, 3 HCP prismatic,
4 HCP pyramidal, 5 BCC110, 6 BCC112. Sequence writes numbered VASP files.
\end{lstlisting}
